%% ============================================================
%% sections/s5-financial.tex
%% H. R. 9510 Bill v5.0 - SEC. 5. FINANCIAL DATA TRANSPARENCY; USER FEES;
%% AUTHORIZATION OF APPROPRIATIONS; BUDGETARY EFFECTS. The operative financial
%% core, NEW in v5.0.
%% Rendered visuals: Table 5 (financial-data elements dictionary), Figure 6
%% (Mermaid sequence, the annual reporting flow, TikZ), Table 6 (authorization
%% of appropriations schedule, R columns), Figure 7 (ASCII money flow), and
%% the boxed budgetary-effects clause (billbox).
%% ============================================================
% \clearpage
\billsec{SEC. 5.}{FINANCIAL DATA TRANSPARENCY; USER FEES; AUTHORIZATION OF
APPROPRIATIONS; BUDGETARY EFFECTS.}\phantomsection\label{toc:s5}

\uscprov{1}{\textbf{(a) Financial Data Transparency.}}
\uscprov{2}{\textbf{(1) Annual Publication.} The Secretary shall publish,
annually and in de-identified aggregate form, the verification cost records
reported under section 515D(k) of the Federal Food, Drug, and Cosmetic Act, in
the manner of the transparency reports under section 1128G of the Social
Security Act (42 U.S.C. 1320a-7h), so that the cost of safety assurance in
Physical AI oncology clinical investigations is public financial data
\cite{usc-open-payments, cfr-403-subpart-i}.}
\uscprov{2}{\textbf{(2) Content and Limits.} Each annual summary shall state,
for each interaction class and fiscal year, the number of verification runs,
the aggregate and per-run direct cost, the share attributable to the
hard-catastrophe-predicate gates, and the human-review cost occasioned by
ESCALATE determinations; and shall contain no patient-identifiable
information and no proprietary line-item information. Table 5 is the
financial-data elements dictionary for this section.}

\clearpage

{\footnotesize
\begin{xltabular}{\textwidth}{@{}L{4.4cm} L{2.3cm} Y L{2.2cm}@{}}
\hline\noalign{\vskip 0.04cm}
\textbf{Financial data element} & \textbf{Who files} & \textbf{Where it flows}
& \textbf{Public form}\\
\hline\noalign{\vskip 0.04cm}
\endfirsthead
\hline\noalign{\vskip 0.04cm}
\textbf{Financial data element} & \textbf{Who files} & \textbf{Where it flows}
& \textbf{Public form}\\
\hline\noalign{\vskip 0.04cm}
\endhead
\hline\noalign{\vskip 0.04cm}
\endlastfoot
Verification cost record, per run (515D(k)(1)) & Sponsor & Hash-chained part 11
audit trail & Not public (audit)\\
\hline\noalign{\vskip 0.04cm}
Fiscal-year aggregate of cost records (515D(k)(3)) & Sponsor & Annual report to
the Secretary & De-identified aggregate\\
\hline\noalign{\vskip 0.04cm}
Annual cost summary by interaction class (SEC. 5(a)) & Secretary & Public
summary, Open Payments manner & Public\\
\hline\noalign{\vskip 0.04cm}
Verification-service financial arrangements (SEC. 5(b)) & Sponsor and
investigator & Part 54 certification, disclosure w/ application & Not
public (review)\\
\hline\noalign{\vskip 0.04cm}
Implementation cost statement (SEC. 5(a)(3)) & Secretary & Annual budget
justification to Congress & Public\\
\hline\noalign{\vskip 0.04cm}
Authorization and outlay series (SEC. 5(d); App. A) & Congress; CBO after
introduction & Appropriations process; scorecards & Public\\
\hline
\end{xltabular}}
\vspace{-0.725cm}
\figcaption{Table 5. The financial-data elements dictionary: each element this
Act creates or touches,\\who files it, where it flows, and its public form.
Figure 6 draws the same flow in time order.}

\uscprov{2}{\textbf{(3) Implementation Cost Statement.} The Secretary shall
include, with the annual budget justification, a statement of the cost of
carrying out section 515D and this section, in the controlled vocabulary of
the Federal budget process \cite{gao-budget-glossary, omb-a11}.}

\begin{mermaidfig}
\node[mmin,text width=20mm]   (sp)  at (0,0)     {Sponsor};
\node[mmstep,text width=20mm] (at)  at (3.0,0)   {Audit trail\\(part 11)};
\node[mmlaw,text width=20mm]  (sec) at (6.0,0)   {Secretary\\(FDA)};
\node[mmstep,text width=20mm] (pub) at (9.0,0)   {Public\\summary};
\node[mmgoal,text width=20mm] (con) at (12.0,0)  {Congress};
\draw[mmedged] (sp.south)  -- ++(0,-5.0);
\draw[mmedged] (at.south)  -- ++(0,-5.0);
\draw[mmedged] (sec.south) -- ++(0,-5.0);
\draw[mmedged] (pub.south) -- ++(0,-5.0);
\draw[mmedged] (con.south) -- ++(0,-5.0);
\draw[mmedge] (0,-1.1)  -- node[mmlabel]{1. cost record per run (515D(k)(1))} (3.0,-1.1);
\draw[mmedge] (0,-2.0)  -- node[mmlabel,pos=0.5]{2. annual fiscal-year report (515D(k)(3))} (6.0,-2.0);
\draw[mmedge] (0,-2.9)  -- node[mmlabel,pos=0.5]{3. part 54-aligned disclosures (SEC. 5(b))} (6.0,-2.9);
\draw[mmedge] (6.0,-3.8) -- node[mmlabel]{4. de-identified aggregate} (9.0,-3.8);
\draw[mmedge] (6.0,-4.7) -- node[mmlabel,pos=0.62]{5. implementation cost statement} (12.0,-4.7);
\end{mermaidfig}
\vspace{-0.6cm}
\figcaption{Figure 6. The annual financial-data reporting flow: the sponsor writes per-run cost records to the
audit trail, reports the fiscal-year aggregate and the part 54-aligned
disclosures to the Secretary, and the Secretary\\publishes the de-identified
summary and delivers the implementation cost statement to Congress.}

\uscprov{1}{\textbf{(b) Financial Disclosure Alignment.} The financial
disclosure discipline of part 54 of title 21, Code of Federal Regulations
(certification on Form FDA 3454 or disclosure on Form FDA 3455), shall apply
to a financial arrangement between a sponsor and a provider of verification
services under section 515D - including compensation affected by the outcome
of a verification run, a proprietary interest in the verification process,
and significant payments of other sorts - to the same extent as if the
provider were a clinical investigator, so the entity that clears the code has
no undisclosed stake in the clearance \cite{cfr-part54}.}

\uscprov{1}{\textbf{(c) User-Fee Treatment.}}
\uscprov{2}{\textbf{(1) No Fee for a Record Alone.} Section 738(a) of the
Federal Food, Drug, and Cosmetic Act (21 U.S.C. 379j(a)) is amended by adding
at the end the new paragraph (4) shown in the comparative print in section 4:
no fee is assessed for the filing of a verification record, or of a
verification cost record, under section 515D, and such a filing, alone, is
not a submission \cite{usc-379j, usc-379i}.}
\uscprov{2}{\textbf{(2) Sense of Congress on Reauthorization.} It is the sense
of Congress that any incremental review workload attributable to verification
records should be priced through the next medical-device user-fee
reauthorization - the successor to the Medical Device User Fee Amendments of
2022, which authorize fees through fiscal year 2027 - rather than a statutory fee \cite{pl-mdufa5, fda-mdufa-rates}.}

\uscprov{1}{\textbf{(d) Authorization of Appropriations.} There are authorized
to be appropriated to the Secretary, for the Food and Drug Administration
Salaries and Expenses account, to carry out section 515D and this section, the
amounts set out in Table 6 - 18,000,000 dollars for fiscal year 2027,
12,000,000 dollars for fiscal year 2028, 10,000,000 dollars for fiscal year
2029, and 9,000,000 dollars for each of fiscal years 2030 and 2031, a total of
58,000,000 dollars (amounts illustrative pending a formal estimate) - to
remain available until expended. Amounts under this subsection are authorized,
not appropriated; budget authority flows only when provided in an
appropriations Act \cite{gao-budget-glossary}. Figure 7 maps the flow.}

{\footnotesize
\begin{xltabular}{\textwidth}{@{}Y R{1.35cm} R{1.35cm} R{1.35cm} R{1.35cm} R{1.35cm} R{1.1cm}@{}}
\hline\noalign{\vskip 0.04cm}
\textbf{Activity (\$ millions, illustrative)} & \textbf{FY27} & \textbf{FY28} &
\textbf{FY29} & \textbf{FY30} & \textbf{FY31} & \textbf{Total}\\
\hline\noalign{\vskip 0.04cm}
\endfirsthead
\hline\noalign{\vskip 0.04cm}
\textbf{Activity (\$ millions, illustrative)} & \textbf{FY27} & \textbf{FY28} &
\textbf{FY29} & \textbf{FY30} & \textbf{FY31} & \textbf{Total}\\
\hline\noalign{\vskip 0.04cm}
\endhead
\hline\noalign{\vskip 0.04cm}
\endlastfoot
Rulemaking under section 515D(h) & 4.0 & 2.0 & 0.5 & 0.5 & 0.5 & 7.5\\
\hline\noalign{\vskip 0.04cm}
Verification-records system (build, operate) & 9.0 & 6.0 & 4.5 & 3.0 &
3.0 & 25.5\\
\hline\noalign{\vskip 0.04cm}
Sequencing inspections and record review & 3.0 & 2.5 & 3.5 & 4.0 & 4.0 &
17.0\\
\hline\noalign{\vskip 0.04cm}
Annual de-identified cost summaries & 2.0 & 1.5 & 1.5 & 1.5 & 1.5 & 8.0\\
\hline\noalign{\vskip 0.04cm}
\textbf{Total authorized} & \textbf{18.0} & \textbf{12.0} & \textbf{10.0} &
\textbf{9.0} & \textbf{9.0} & \textbf{58.0}\\
\hline
\end{xltabular}}
\vspace{-0.5cm}
\figcaption{Table 6. The authorization of appropriations schedule by activity
and fiscal year (dollars\\in millions, illustrative pending a formal estimate;
right-aligned R columns). Every\\column and row sums; Appendix A converts this
budget authority to estimated outlays.}

\begin{asciifig}
            HOW SECTION 515D IMPLEMENTATION IS PAID FOR (SEC. 5)

Congress: annual appropriations          Industry: MDUFA user fees
(authorized by SEC. 5(d))                (21 U.S.C. 379i-379j, unchanged)
             |                                        |
             |  budget authority                      |  fee revenue
             v                                        v
+-----------------------------------------------------------------------+
|                     FDA device program (S&E account)                  |
+-----------------------------------------------------------------------+
    |                    |                    |                 |
    v                    v                    v                 v
rulemaking         verification          sequencing       annual public
515D(h)            records (FY27-28      inspections      cost summaries
(FY27 heavy)       build, then O&M)      (steady state)   (de-identified)
                      
NOT permitted: a new statutory fee for filing a verification record
(SEC. 5(c)); pricing belongs to the MDUFA VI negotiation (by 10/2027)
\end{asciifig}
\vspace{-0.4cm}
\figcaption{Figure 7. The money flow (ASCII): the two funding sources, the
account they meet in,\\the four activities of Table 6 they pay for, and the
fee treatment boundary SEC. 5(c) draws.}

\uscprov{1}{\textbf{(e) Budgetary Effects.} This Act creates no direct
spending and no revenue effect; its costs are subject to appropriations. For
scorekeeping, the Act carries the standard clause \cite{usc-paygo-2010,
house-rule-xxi}:}

\billbox{\textbf{Budgetary Effects.} The budgetary effects of this Act, for
the purpose of complying with the Statutory Pay-As-You-Go Act of 2010, shall
be determined by reference to the latest statement titled ``Budgetary Effects
of PAYGO Legislation'' for this Act, submitted for printing in the
Congressional Record by the Chairman of the House Committee on the Budget,
provided that such statement has been submitted prior to the vote on passage.}
